\subsection{Observability-krav (O1.1--O1.3)}
\label{sec:eval-observability}

Krav~O1.1--O1.3 ble i hovedsak oppfylt, men arbeidet med observability
ble mer krevende enn først antatt. Den største lærdommen var at
utfordringen ikke først og fremst lå i å lage logger, helsesjekker eller
metrikker, men i å få Prometheus, Grafana og API-et til å fungere stabilt
sammen i et containerisert miljø.

Det mest tidkrevende problemet var å få Prometheus til å scrape
\texttt{/metrics}-endepunktet fra API-containeren. I teorien er dette bare
en HTTP-forespørsel, men i praksis støtte den på flere lag i
applikasjonen og infrastrukturen. Intern Docker-DNS, \texttt{AllowedHosts},
\texttt{RequireHost} og HTTPS-redirect påvirket alle om Prometheus faktisk
fikk svar fra API-et. Vi prøvde flere varianter før vi landet på en løsning
der Prometheus scraper over HTTP internt i Docker-nettverket, samtidig som
\texttt{/metrics} eksplisitt unntas fra HTTPS-redirect. Dette var et
praktisk kompromiss: TLS gir liten verdi internt mellom containere på samme
host, men HTTPS bør fortsatt brukes for ekstern trafikk.

Grafana ga også noen uventede problemer. Datasource-oppsettet fungerte
ikke stabilt før UID-en til Prometheus-datasourcen ble hardkodet i
dashboard-filene. Dette er ikke den mest fleksible løsningen, men den gjorde
provisioning mer forutsigbar. I tillegg måtte vi håndtere problemer med
plugin-oppdateringer og minnebruk i Grafana-containeren. Dette viste at
observability-verktøyene selv også må driftes og konfigureres riktig, ikke
bare legges til som en ekstra tjeneste i Docker Compose.

Et annet konkret eksempel var Docker-healthchecken. Den brukte
\texttt{curl}, men base-imaget inneholdt ikke \texttt{curl}. Resultatet var
at containeren ble markert som \texttt{unhealthy}, selv om API-et i seg selv
kunne kjøre. Feilen var enkel å rette, men vanskeligere å finne fordi
symptomet oppstod i samspillet mellom Docker, API-containeren og
Prometheus. Dette gjorde det tydelig at flere av problemene var
diagnostiseringsproblemer mer enn rene kodeproblemer.

En viktig erfaring er derfor at observability ikke bare handler om å
eksponere riktige data, men også om å kunne stole på infrastrukturen som
samler dem inn. Flere av løsningene ble pragmatiske kompromisser, som HTTP
internt i Docker-nettet, hardkodet datasource-UID og eksponering av Grafana
og Prometheus-porter for at HDO lettere kunne teste. Disse valgene var
nyttige i prosjektfasen, men bør vurderes strengere før produksjonsbruk.

Den viktigste svakheten i dagens løsning er at \texttt{/metrics}-endepunktet
er eksponert uten autentisering. Før produksjon bør dette sikres med
nettverksbegrensning, reverse proxy-regler eller Bearer-autentisering. I
tillegg burde testdekningen vært høyere, og feilsøkingsloggen for
observability-oppsettet mer systematisk, slik at slike problemer kunne vært
oppdaget tidligere.